% OC Core 1.3.3 -- Cross-level structures public chapter

\section{Cross-Level Structures}
\label{sec:crossk_master}

\subsection{Purpose of the Cross-K layer}
The Cross-K layer formalizes relations between declared continuum levels. Each
level has its own axes, potentials, thresholds, flows, cycles, domains,
boundaries, and continuumness functionals; many scientific questions become
visible only when two adjacent levels must be compared. The Cross-K layer is
therefore a scientific layer, not a build module: it asks which structure is
retained, which observable is added, which reduction is attempted, and when a
lawful demotion is allowed.

\subsection{Definition of Cross-K structures}
A Cross-K structure relates at least two levels \(K_i\) and \(K_j\) with
\(i \ne j\). It records a relation \(R_{ij}\) between the components of the
corresponding continua and states the compatibility condition under which a
transition, projection, lift, boundary transfer, or cycle comparison is valid.
For OC Core 1.3.3, this language is bounded by the public K-level theorem and
finite semantic witness suite: it is a declared witness taxonomy, not an
unrestricted proof of every possible higher-order hierarchy.

\subsection{Reader-facing transition logic}
The detailed Cross-K chapters that follow this one provide the transition
content. They should be read in three passes. First, identify the added
observable or structural axis. Second, ask whether a lower-level reduction keeps
the relevant witness. Third, check whether the added observable is inert; if it
is inert, lawful demotion is permitted, and if it is retained, demotion is
blocked by the declared verdict suite.

\subsection{Boundary of the public claim}
K0 through K12 are retained as a declared witness taxonomy under stated
assumptions. The public release does not claim that K11 or K12 complete all
future metatheory, bounded cross-domain science, or scoped comparative unrestricted comparative claim.
They mark upper-level witness obligations and research-taxonomy candidates
whose promoted force depends on the theorem sheets, finite cases, comparator
rows, and reopening rules included in the release evidence package.
